%% Generated by Sphinx.
\def\sphinxdocclass{report}
\documentclass[letterpaper,10pt,english]{sphinxmanual}
\ifdefined\pdfpxdimen
   \let\sphinxpxdimen\pdfpxdimen\else\newdimen\sphinxpxdimen
\fi \sphinxpxdimen=.75bp\relax
\ifdefined\pdfimageresolution
    \pdfimageresolution= \numexpr \dimexpr1in\relax/\sphinxpxdimen\relax
\fi
%% let collapsible pdf bookmarks panel have high depth per default
\PassOptionsToPackage{bookmarksdepth=5}{hyperref}

\PassOptionsToPackage{booktabs}{sphinx}
\PassOptionsToPackage{colorrows}{sphinx}

\PassOptionsToPackage{warn}{textcomp}
\usepackage[utf8]{inputenc}
\ifdefined\DeclareUnicodeCharacter
% support both utf8 and utf8x syntaxes
  \ifdefined\DeclareUnicodeCharacterAsOptional
    \def\sphinxDUC#1{\DeclareUnicodeCharacter{"#1}}
  \else
    \let\sphinxDUC\DeclareUnicodeCharacter
  \fi
  \sphinxDUC{00A0}{\nobreakspace}
  \sphinxDUC{2500}{\sphinxunichar{2500}}
  \sphinxDUC{2502}{\sphinxunichar{2502}}
  \sphinxDUC{2514}{\sphinxunichar{2514}}
  \sphinxDUC{251C}{\sphinxunichar{251C}}
  \sphinxDUC{2572}{\textbackslash}
\fi
\usepackage{cmap}
\usepackage[T1]{fontenc}
\usepackage{amsmath,amssymb,amstext}
\usepackage{babel}



\usepackage{tgtermes}
\usepackage{tgheros}
\renewcommand{\ttdefault}{txtt}



\usepackage[Bjarne]{fncychap}
\usepackage{sphinx}

\fvset{fontsize=auto}
\usepackage{geometry}


% Include hyperref last.
\usepackage{hyperref}
% Fix anchor placement for figures with captions.
\usepackage{hypcap}% it must be loaded after hyperref.
% Set up styles of URL: it should be placed after hyperref.
\urlstyle{same}

\addto\captionsenglish{\renewcommand{\contentsname}{Contents:}}

\usepackage{sphinxmessages}
\setcounter{tocdepth}{1}



\title{chowclassifier}
\date{Mar 29, 2023}
\release{1.0.5}
\author{Luc T Miaz}
\newcommand{\sphinxlogo}{\vbox{}}
\renewcommand{\releasename}{Release}
\makeindex
\begin{document}

\ifdefined\shorthandoff
  \ifnum\catcode`\=\string=\active\shorthandoff{=}\fi
  \ifnum\catcode`\"=\active\shorthandoff{"}\fi
\fi

\pagestyle{empty}
\sphinxmaketitle
\pagestyle{plain}
\sphinxtableofcontents
\pagestyle{normal}
\phantomsection\label{\detokenize{index::doc}}


\begin{sphinxShadowBox}
\sphinxstyletopictitle{Table of Contents}
\begin{itemize}
\item {} 
\sphinxAtStartPar
\phantomsection\label{\detokenize{index:id4}}{\hyperref[\detokenize{index:welcome-to-chowclassifier-s-documentation}]{\sphinxcrossref{Welcome to chowclassifier’s documentation!}}}

\item {} 
\sphinxAtStartPar
\phantomsection\label{\detokenize{index:id5}}{\hyperref[\detokenize{index:module-description}]{\sphinxcrossref{Module Description}}}
\begin{itemize}
\item {} 
\sphinxAtStartPar
\phantomsection\label{\detokenize{index:id6}}{\hyperref[\detokenize{index:introduction}]{\sphinxcrossref{Introduction}}}

\item {} 
\sphinxAtStartPar
\phantomsection\label{\detokenize{index:id7}}{\hyperref[\detokenize{index:chow-test}]{\sphinxcrossref{Chow test}}}

\item {} 
\sphinxAtStartPar
\phantomsection\label{\detokenize{index:id8}}{\hyperref[\detokenize{index:installation}]{\sphinxcrossref{Installation}}}

\item {} 
\sphinxAtStartPar
\phantomsection\label{\detokenize{index:id9}}{\hyperref[\detokenize{index:use}]{\sphinxcrossref{Use}}}

\end{itemize}

\item {} 
\sphinxAtStartPar
\phantomsection\label{\detokenize{index:id10}}{\hyperref[\detokenize{index:class-and-method-descriptions}]{\sphinxcrossref{Class and method descriptions}}}

\item {} 
\sphinxAtStartPar
\phantomsection\label{\detokenize{index:id11}}{\hyperref[\detokenize{index:licence}]{\sphinxcrossref{Licence}}}

\item {} 
\sphinxAtStartPar
\phantomsection\label{\detokenize{index:id12}}{\hyperref[\detokenize{index:outputs}]{\sphinxcrossref{Outputs}}}
\begin{itemize}
\item {} 
\sphinxAtStartPar
\phantomsection\label{\detokenize{index:id13}}{\hyperref[\detokenize{index:plot-and-plot-individually}]{\sphinxcrossref{\sphinxcode{\sphinxupquote{plot}} and \sphinxcode{\sphinxupquote{plot\_individually}}}}}

\item {} 
\sphinxAtStartPar
\phantomsection\label{\detokenize{index:id14}}{\hyperref[\detokenize{index:save-to-file}]{\sphinxcrossref{\sphinxcode{\sphinxupquote{save\_to\_file}}}}}

\end{itemize}

\item {} 
\sphinxAtStartPar
\phantomsection\label{\detokenize{index:id15}}{\hyperref[\detokenize{index:version-changes}]{\sphinxcrossref{Version changes}}}

\end{itemize}
\end{sphinxShadowBox}


\chapter{Module Description}
\label{\detokenize{index:module-description}}

\section{Introduction}
\label{\detokenize{index:introduction}}
\sphinxAtStartPar
There are many application where it is useful to analyse rapidly a large number of time\sphinxhyphen{}series for overall trends and possible breakpoint at which the trend changes. For example, spotting increasing trends in potential contaminants concentration from the outputs of non\sphinxhyphen{}target screening from high resolution mass spectrometry.

\sphinxAtStartPar
ChowClassifier aims to solve this. It is a script that performs a series of Chow tests on a time\sphinxhyphen{}indexed dataset to determine if a breakpoint in the time trends is probable. Two classes \sphinxtitleref{Chow} and \sphinxtitleref{ChewData} enable to perform the tests, generate plots and export the results in \sphinxtitleref{csv} and \sphinxtitleref{excel}. Dataset with groups can be ran simultaneously with \sphinxtitleref{ChewData} indicating with \sphinxtitleref{namecol} the name of the grouping column. Calling \sphinxtitleref{run} on a \sphinxtitleref{ChewData} instance will create an instance of \sphinxtitleref{Chow} for each group and classify them into the following categories.


\begin{savenotes}\sphinxattablestart
\sphinxthistablewithglobalstyle
\centering
\begin{tabulary}{\linewidth}[t]{TTT}
\sphinxtoprule
\sphinxtableatstartofbodyhook
\sphinxAtStartPar
Breakpoint ?
&
\sphinxAtStartPar
Category
&
\sphinxAtStartPar
Description
\\
\sphinxhline
\sphinxAtStartPar
No
&
\sphinxAtStartPar
N
&
\sphinxAtStartPar
non\sphinxhyphen{}significant overall trend
\\
\sphinxhline
\sphinxAtStartPar
No
&
\sphinxAtStartPar
I
&
\sphinxAtStartPar
significant increasing overall trend
\\
\sphinxhline
\sphinxAtStartPar
No
&
\sphinxAtStartPar
D
&
\sphinxAtStartPar
significant decreasing overall trend
\\
\sphinxhline
\sphinxAtStartPar
Yes
&
\sphinxAtStartPar
NN
&
\sphinxAtStartPar
non\sphinxhyphen{}significant trend on set1 and non\sphinxhyphen{}significant trend on set2
\\
\sphinxhline
\sphinxAtStartPar
Yes
&
\sphinxAtStartPar
NI
&
\sphinxAtStartPar
non\sphinxhyphen{}significant trend on set1 and significant increasing trend on set2
\\
\sphinxhline
\sphinxAtStartPar
Yes
&
\sphinxAtStartPar
ND
&
\sphinxAtStartPar
non\sphinxhyphen{}significant trend on set1 ans significant decreasing trend on set2
\\
\sphinxhline
\sphinxAtStartPar
Yes
&
\sphinxAtStartPar
IN
&
\sphinxAtStartPar
significant increasing trend on set1 and non\sphinxhyphen{}significant trend on set2
\\
\sphinxhline
\sphinxAtStartPar
Yes
&
\sphinxAtStartPar
ID
&
\sphinxAtStartPar
significant increasing trend on set1 and significant decreasing trend on set2
\\
\sphinxhline
\sphinxAtStartPar
Yes
&
\sphinxAtStartPar
iI
&
\sphinxAtStartPar
significant increasing trend on both set1 and set2 with greater increase in set2
\\
\sphinxhline
\sphinxAtStartPar
Yes
&
\sphinxAtStartPar
Ii
&
\sphinxAtStartPar
significant increasing trend on both set1 and set2 with greater increase in set1
\\
\sphinxhline
\sphinxAtStartPar
Yes
&
\sphinxAtStartPar
DN
&
\sphinxAtStartPar
significant decreasing trend on set1 and non\sphinxhyphen{}significant trend on set2
\\
\sphinxhline
\sphinxAtStartPar
Yes
&
\sphinxAtStartPar
DI
&
\sphinxAtStartPar
significant decreasing trend on set1 and significant increasing trend on set2
\\
\sphinxhline
\sphinxAtStartPar
Yes
&
\sphinxAtStartPar
dD
&
\sphinxAtStartPar
significant decreasing trend on both set1 and set2 with greater decrease in set2
\\
\sphinxhline
\sphinxAtStartPar
Yes
&
\sphinxAtStartPar
Dd
&
\sphinxAtStartPar
significant decreasing trend on both set1 and set2 with greater decrease in set1
\\
\sphinxbottomrule
\end{tabulary}
\sphinxtableafterendhook\par
\sphinxattableend\end{savenotes}

\sphinxAtStartPar
where set1 and set2 indicate respectively the data before the breakpoint and after. The decision process for this classification is shown in the following schema:

\noindent{\hspace*{\fill}\sphinxincludegraphics{{schema}.png}}


\section{Chow test}
\label{\detokenize{index:chow-test}}
\sphinxAtStartPar
Chow test was first derived by Gregory Chow %
\begin{footnote}[1]\sphinxAtStartFootnote
Chow, Gregory C. “Test of Equality Between Sets of Coefficients in Two Linear Regressions”. \sphinxstyleemphasis{Econometrica} 28, no. 3 (1960): 591\sphinxhyphen{}605. \sphinxurl{https://www.jstor.org/stable/1910133}.
%
\end{footnote} in 1960 and later by Franklin Fisher %
\begin{footnote}[2]\sphinxAtStartFootnote
Fisher, Franklin M. “Tests of Equality Between Sets of Coefficients in Two Linear Regressions: An Expository Note”. \sphinxstyleemphasis{Econometrica2} 38, no. 2 (1970): 361\sphinxhyphen{}66. \sphinxurl{https://www.jstor.org/stable/1913018}.
%
\end{footnote}. We test the significance of the breakpoint at level \(\alpha^\star = \frac{\alpha}{m}\) (where \(m\) is the number of breakpoints tests, thus \(\alpha^\star\) is the Bonferroni corrected significance level)  under the null hypothesis \(Z=\frac{S_c-(S_1+S_2)}{S_1+S_2}\cdot\frac{N_1+N_2-2\cdot k}{k}\) where \(k=3\) is the total number of parameters. This follows a F\sphinxhyphen{}distribution with \(k\) and \(N_1+N_2-6\) degrees of freedom, where \(S_C`\) is the sum of squared residuals of the regression on the full time series, \(S_1\), \(S_2\) are the sums of squared residuals of the regression on the first, and respectively, the second half of the time series and \(N_1\), \(N_2\) are the number of observation for each half. See also %
\begin{footnote}[3]\sphinxAtStartFootnote
\sphinxurl{https://en.wikipedia.org/wiki/Chow\_test}
%
\end{footnote}.
\subsubsection*{References}


\section{Installation}
\label{\detokenize{index:installation}}
\sphinxAtStartPar
Use the following command to install chowlcassifier:

\begin{sphinxVerbatim}[commandchars=\\\{\}]
\PYG{n}{pip} \PYG{n}{install} \PYG{n}{chowclassifier}
\end{sphinxVerbatim}


\section{Use}
\label{\detokenize{index:use}}
\sphinxAtStartPar
The script can be ran on any file with \sphinxcode{\sphinxupquote{python \sphinxhyphen{}m chowclassifier \sphinxhyphen{}f path/to/file/filename.csv \sphinxhyphen{}X xcol \sphinxhyphen{}y ycol \sphinxhyphen{}n grouping\_name}} or it can be imported:

\begin{sphinxVerbatim}[commandchars=\\\{\}]
\PYG{k+kn}{from} \PYG{n+nn}{chowclassifier} \PYG{k+kn}{import} \PYG{n}{ChewData}
\PYG{c+c1}{\PYGZsh{} define the path to data}
\PYG{n}{filepath} \PYG{o}{=} \PYG{l+s+s2}{\PYGZdq{}}\PYG{l+s+s2}{example/data}\PYG{l+s+s2}{\PYGZdq{}}
\PYG{c+c1}{\PYGZsh{} define the name of the file containing the data (with extension)}
\PYG{n}{filename} \PYG{o}{=} \PYG{l+s+s2}{\PYGZdq{}}\PYG{l+s+s2}{stocks.csv}\PYG{l+s+s2}{\PYGZdq{}}
\PYG{c+c1}{\PYGZsh{} define the path where figures will be saved}
\PYG{n}{savingpath\PYGZus{}figures} \PYG{o}{=} \PYG{l+s+s1}{\PYGZsq{}}\PYG{l+s+s1}{example/fig}\PYG{l+s+s1}{\PYGZsq{}}
\PYG{c+c1}{\PYGZsh{} what name has the X/time column?}
\PYG{n}{timecol} \PYG{o}{=} \PYG{l+s+s1}{\PYGZsq{}}\PYG{l+s+s1}{Date}\PYG{l+s+s1}{\PYGZsq{}}
\PYG{c+c1}{\PYGZsh{} what name has the y/value column?}
\PYG{n}{ycol} \PYG{o}{=} \PYG{l+s+s1}{\PYGZsq{}}\PYG{l+s+s1}{Close}\PYG{l+s+s1}{\PYGZsq{}}
\PYG{c+c1}{\PYGZsh{} what level of confidence? Note, when}
\PYG{c+c1}{\PYGZsh{} testing multiple breakpoint, a Bonferroni correction will be applied}
\PYG{n}{alpha} \PYG{o}{=} \PYG{l+m+mf}{0.01}
\PYG{c+c1}{\PYGZsh{} initial breakpoint ?}
\PYG{n}{initial\PYGZus{}breakpoint} \PYG{o}{=} \PYG{k+kc}{None}

\PYG{n}{C} \PYG{o}{=} \PYG{n}{ChewData}\PYG{p}{(}\PYG{n}{filename} \PYG{o}{=} \PYG{l+s+sa}{f}\PYG{l+s+s2}{\PYGZdq{}}\PYG{l+s+si}{\PYGZob{}}\PYG{l+s+s1}{\PYGZsq{}}\PYG{l+s+s1}{/}\PYG{l+s+s1}{\PYGZsq{}}\PYG{o}{.}\PYG{n}{join}\PYG{p}{(}\PYG{p}{[}\PYG{n}{x} \PYG{k}{for} \PYG{n}{x} \PYG{o+ow}{in} \PYG{p}{[}\PYG{n}{filepath}\PYG{p}{,}\PYG{n}{filename}\PYG{p}{]} \PYG{k}{if} \PYG{n}{x} \PYG{o+ow}{not} \PYG{o+ow}{in} \PYG{p}{[}\PYG{k+kc}{None}\PYG{p}{,}\PYG{l+s+s1}{\PYGZsq{}}\PYG{l+s+s1}{\PYGZsq{}}\PYG{p}{]}\PYG{p}{]}\PYG{p}{)}\PYG{l+s+si}{\PYGZcb{}}\PYG{l+s+s2}{\PYGZdq{}}\PYG{p}{,}
            \PYG{n}{timecol} \PYG{o}{=} \PYG{n}{timecol}\PYG{p}{,}\PYG{c+c1}{\PYGZsh{} name of time column (used as x\PYGZhy{}axis)}
            \PYG{n}{ycol} \PYG{o}{=} \PYG{n}{ycol}\PYG{p}{,}\PYG{c+c1}{\PYGZsh{} name for value column (used as y\PYGZhy{}axis) Leave blank if multiple}
            \PYG{n}{alpha} \PYG{o}{=} \PYG{n}{alpha}\PYG{c+c1}{\PYGZsh{} level of confidence to use}
            \PYG{p}{)}

\PYG{c+c1}{\PYGZsh{}\PYGZsh{}\PYGZsh{}\PYGZsh{}\PYGZsh{}\PYGZsh{}\PYGZsh{}\PYGZsh{}\PYGZsh{}\PYGZsh{} use the code line below if your X/time column is not a number, e.g. a date}
\PYG{c+c1}{\PYGZsh{}\PYGZsh{}\PYGZsh{}\PYGZsh{}\PYGZsh{}\PYGZsh{}\PYGZsh{}\PYGZsh{}\PYGZsh{}\PYGZsh{} comment it otherwise}
\PYG{c+c1}{\PYGZsh{} C.parse\PYGZus{}timecol(date\PYGZus{}format = \PYGZsq{}\PYGZpc{}Y\PYGZhy{}\PYGZpc{}m\PYGZhy{}\PYGZpc{}d\PYGZsq{}) \PYGZsh{}}

\PYG{n}{C}\PYG{o}{.}\PYG{n}{run}\PYG{p}{(}\PYG{n}{initial\PYGZus{}breakpoint} \PYG{o}{=} \PYG{n}{initial\PYGZus{}breakpoint}\PYG{p}{)}
\PYG{n}{C}\PYG{o}{.}\PYG{n}{plot}\PYG{p}{(}\PYG{n}{xlabel} \PYG{o}{=} \PYG{l+s+s1}{\PYGZsq{}}\PYG{l+s+s1}{time}\PYG{l+s+s1}{\PYGZsq{}}\PYG{p}{,}\PYG{c+c1}{\PYGZsh{} label for x axis}
      \PYG{n}{ylabel}\PYG{o}{=}\PYG{l+s+s1}{\PYGZsq{}}\PYG{l+s+s1}{value}\PYG{l+s+s1}{\PYGZsq{}}\PYG{p}{,}\PYG{c+c1}{\PYGZsh{} label for y axis}
      \PYG{n}{title} \PYG{o}{=} \PYG{l+s+s2}{\PYGZdq{}}\PYG{l+s+s2}{Chow Classification}\PYG{l+s+s2}{\PYGZdq{}}\PYG{p}{,}\PYG{c+c1}{\PYGZsh{} title of the main plot}
      \PYG{n}{filename} \PYG{o}{=} \PYG{l+s+sa}{f}\PYG{l+s+s2}{\PYGZdq{}}\PYG{l+s+si}{\PYGZob{}}\PYG{l+s+s1}{\PYGZsq{}}\PYG{l+s+s1}{/}\PYG{l+s+s1}{\PYGZsq{}}\PYG{o}{.}\PYG{n}{join}\PYG{p}{(}\PYG{p}{[}\PYG{n}{x} \PYG{k}{for} \PYG{n}{x} \PYG{o+ow}{in} \PYG{p}{[}\PYG{n}{savingpath\PYGZus{}figures}\PYG{p}{,}\PYG{n}{filename}\PYG{p}{]} \PYG{k}{if} \PYG{n}{x} \PYG{o+ow}{not} \PYG{o+ow}{in} \PYG{p}{[}\PYG{k+kc}{None}\PYG{p}{,}\PYG{l+s+s1}{\PYGZsq{}}\PYG{l+s+s1}{\PYGZsq{}}\PYG{p}{]}\PYG{p}{]}\PYG{p}{)}\PYG{l+s+si}{\PYGZcb{}}\PYG{l+s+s2}{.png}\PYG{l+s+s2}{\PYGZdq{}}\PYG{p}{,}\PYG{c+c1}{\PYGZsh{} filename for the figure}
      \PYG{n}{figsize}\PYG{o}{=}\PYG{p}{(}\PYG{l+m+mi}{16}\PYG{p}{,}\PYG{l+m+mi}{8}\PYG{p}{)}\PYG{p}{,}\PYG{c+c1}{\PYGZsh{} figure size}
      \PYG{n}{sharey}\PYG{o}{=}\PYG{k+kc}{True} \PYG{c+c1}{\PYGZsh{} whether the individual plots are forced to share y axis}
      \PYG{p}{)}
\PYG{n}{C}\PYG{o}{.}\PYG{n}{plot\PYGZus{}individually}\PYG{p}{(}\PYG{n}{savingpath} \PYG{o}{=} \PYG{n}{savingpath\PYGZus{}figures}\PYG{p}{,}\PYG{c+c1}{\PYGZsh{} Saving path for the figures}
                  \PYG{n+nb}{format} \PYG{o}{=} \PYG{l+s+s1}{\PYGZsq{}}\PYG{l+s+s1}{png}\PYG{l+s+s1}{\PYGZsq{}}\PYG{p}{,}\PYG{c+c1}{\PYGZsh{} figure format}
                  \PYG{n}{xlabel} \PYG{o}{=} \PYG{l+s+s1}{\PYGZsq{}}\PYG{l+s+s1}{time}\PYG{l+s+s1}{\PYGZsq{}}\PYG{p}{,}\PYG{c+c1}{\PYGZsh{} label for x axis}
                  \PYG{n}{ylabel}\PYG{o}{=}\PYG{l+s+s1}{\PYGZsq{}}\PYG{l+s+s1}{value}\PYG{l+s+s1}{\PYGZsq{}}\PYG{p}{,}\PYG{c+c1}{\PYGZsh{} label for y axis}
                  \PYG{p}{)}

\PYG{n}{C}\PYG{o}{.}\PYG{n}{plot\PYGZus{}by\PYGZus{}group}\PYG{p}{(}\PYG{l+s+s1}{\PYGZsq{}}\PYG{l+s+s1}{g}\PYG{l+s+s1}{\PYGZsq{}}\PYG{p}{,}
               \PYG{n}{savingpath} \PYG{o}{=} \PYG{n}{savingpath\PYGZus{}figures}\PYG{p}{,}\PYG{c+c1}{\PYGZsh{} Saving path for the figures}
               \PYG{n+nb}{format} \PYG{o}{=} \PYG{l+s+s1}{\PYGZsq{}}\PYG{l+s+s1}{png}\PYG{l+s+s1}{\PYGZsq{}}\PYG{p}{,}\PYG{c+c1}{\PYGZsh{} figure format}
               \PYG{n}{xlabel} \PYG{o}{=} \PYG{l+s+s1}{\PYGZsq{}}\PYG{l+s+s1}{time}\PYG{l+s+s1}{\PYGZsq{}}\PYG{p}{,}\PYG{c+c1}{\PYGZsh{} label for x axis}
               \PYG{n}{ylabel}\PYG{o}{=}\PYG{l+s+s1}{\PYGZsq{}}\PYG{l+s+s1}{value}\PYG{l+s+s1}{\PYGZsq{}}\PYG{p}{,}\PYG{c+c1}{\PYGZsh{} label for y axis}
               \PYG{n}{plot\PYGZus{}overall} \PYG{o}{=} \PYG{k+kc}{True}\PYG{p}{,}\PYG{c+c1}{\PYGZsh{} Whether to plot an overall trend across groups (including confidence inverval fill)}
               \PYG{n}{plot\PYGZus{}individual\PYGZus{}fill} \PYG{o}{=} \PYG{k+kc}{True}\PYG{c+c1}{\PYGZsh{} Whether to plot individual confidence inverval fills}
               \PYG{p}{)}
\end{sphinxVerbatim}


\chapter{Class and method descriptions}
\label{\detokenize{index:class-and-method-descriptions}}
\sphinxstepscope


\section{chowclassifier}
\label{\detokenize{source/modules:chowclassifier}}\label{\detokenize{source/modules::doc}}
\sphinxstepscope


\subsection{chowclassifier package}
\label{\detokenize{source/chowclassifier:chowclassifier-package}}\label{\detokenize{source/chowclassifier::doc}}

\subsubsection{Submodules}
\label{\detokenize{source/chowclassifier:submodules}}

\subsubsection{chowclassifier.chow module}
\label{\detokenize{source/chowclassifier:module-chowclassifier.chow}}\label{\detokenize{source/chowclassifier:chowclassifier-chow-module}}\index{module@\spxentry{module}!chowclassifier.chow@\spxentry{chowclassifier.chow}}\index{chowclassifier.chow@\spxentry{chowclassifier.chow}!module@\spxentry{module}}
\sphinxAtStartPar
Implementation of a Chow Classifier
\index{ChewData (class in chowclassifier.chow)@\spxentry{ChewData}\spxextra{class in chowclassifier.chow}}

\begin{fulllineitems}
\phantomsection\label{\detokenize{source/chowclassifier:chowclassifier.chow.ChewData}}
\pysigstartsignatures
\pysiglinewithargsret{\sphinxbfcode{\sphinxupquote{class\DUrole{w}{  }}}\sphinxcode{\sphinxupquote{chowclassifier.chow.}}\sphinxbfcode{\sphinxupquote{ChewData}}}{\emph{\DUrole{n}{df}\DUrole{p}{:}\DUrole{w}{  }\DUrole{n}{DataFrame\DUrole{w}{  }\DUrole{p}{|}\DUrole{w}{  }None}\DUrole{w}{  }\DUrole{o}{=}\DUrole{w}{  }\DUrole{default_value}{None}}, \emph{\DUrole{n}{filename}\DUrole{p}{:}\DUrole{w}{  }\DUrole{n}{str\DUrole{w}{  }\DUrole{p}{|}\DUrole{w}{  }None}\DUrole{w}{  }\DUrole{o}{=}\DUrole{w}{  }\DUrole{default_value}{None}}, \emph{\DUrole{n}{namecol}\DUrole{p}{:}\DUrole{w}{  }\DUrole{n}{str}\DUrole{w}{  }\DUrole{o}{=}\DUrole{w}{  }\DUrole{default_value}{\textquotesingle{}name\textquotesingle{}}}, \emph{\DUrole{n}{initial\_breakpoint}\DUrole{p}{:}\DUrole{w}{  }\DUrole{n}{float\DUrole{w}{  }\DUrole{p}{|}\DUrole{w}{  }None}\DUrole{w}{  }\DUrole{o}{=}\DUrole{w}{  }\DUrole{default_value}{None}}, \emph{\DUrole{n}{breakpoints}\DUrole{p}{:}\DUrole{w}{  }\DUrole{n}{list\DUrole{w}{  }\DUrole{p}{|}\DUrole{w}{  }None}\DUrole{w}{  }\DUrole{o}{=}\DUrole{w}{  }\DUrole{default_value}{None}}, \emph{\DUrole{n}{sep}\DUrole{o}{=}\DUrole{default_value}{\textquotesingle{},\textquotesingle{}}}, \emph{\DUrole{n}{timecol}\DUrole{p}{:}\DUrole{w}{  }\DUrole{n}{str}\DUrole{w}{  }\DUrole{o}{=}\DUrole{w}{  }\DUrole{default_value}{\textquotesingle{}year\textquotesingle{}}}, \emph{\DUrole{n}{ycol}\DUrole{p}{:}\DUrole{w}{  }\DUrole{n}{str}\DUrole{w}{  }\DUrole{o}{=}\DUrole{w}{  }\DUrole{default_value}{\textquotesingle{}value\textquotesingle{}}}, \emph{\DUrole{n}{margin}\DUrole{p}{:}\DUrole{w}{  }\DUrole{n}{int}\DUrole{w}{  }\DUrole{o}{=}\DUrole{w}{  }\DUrole{default_value}{2}}, \emph{\DUrole{n}{alpha}\DUrole{p}{:}\DUrole{w}{  }\DUrole{n}{float}\DUrole{w}{  }\DUrole{o}{=}\DUrole{w}{  }\DUrole{default_value}{0.01}}}{}
\pysigstopsignatures
\sphinxAtStartPar
Bases: \sphinxcode{\sphinxupquote{object}}

\sphinxAtStartPar
Handle a full data frame with grouping of time series
\begin{quote}\begin{description}
\sphinxlineitem{Params df}
\sphinxAtStartPar
dataset (as a pandas DataFrame) or

\sphinxlineitem{Params filename}
\sphinxAtStartPar
path to dataframe (used if df not given)

\sphinxlineitem{Params namecol}
\sphinxAtStartPar
name of the column with group names (used in groupby)

\sphinxlineitem{Params initial\_breakpoint}
\sphinxAtStartPar
initial breakpoint

\sphinxlineitem{Params breakpoints}
\sphinxAtStartPar
candidates for breakpoints

\sphinxlineitem{Params sep}
\sphinxAtStartPar
separator in csv

\sphinxlineitem{Params timecol}
\sphinxAtStartPar
name of the time column in dataset

\sphinxlineitem{Params ycol}
\sphinxAtStartPar
name of the variable column in dataset

\sphinxlineitem{Params margin}
\sphinxAtStartPar
range of indexes around breakpoint/midpoint, see Chow.get\_breakpoints

\sphinxlineitem{Params alpha}
\sphinxAtStartPar
significance level at which to test the trends (Bonferroni correction will be applied later)

\end{description}\end{quote}
\index{get\_df() (chowclassifier.chow.ChewData method)@\spxentry{get\_df()}\spxextra{chowclassifier.chow.ChewData method}}

\begin{fulllineitems}
\phantomsection\label{\detokenize{source/chowclassifier:chowclassifier.chow.ChewData.get_df}}
\pysigstartsignatures
\pysiglinewithargsret{\sphinxbfcode{\sphinxupquote{get\_df}}}{\emph{\DUrole{n}{filename}}, \emph{\DUrole{n}{sep}}}{}
\pysigstopsignatures
\sphinxAtStartPar
get dataframe from file

\end{fulllineitems}

\index{melt() (chowclassifier.chow.ChewData method)@\spxentry{melt()}\spxextra{chowclassifier.chow.ChewData method}}

\begin{fulllineitems}
\phantomsection\label{\detokenize{source/chowclassifier:chowclassifier.chow.ChewData.melt}}
\pysigstartsignatures
\pysiglinewithargsret{\sphinxbfcode{\sphinxupquote{melt}}}{}{}
\pysigstopsignatures
\sphinxAtStartPar
pivot the dataframe from wide to long

\end{fulllineitems}

\index{parse\_timecol() (chowclassifier.chow.ChewData method)@\spxentry{parse\_timecol()}\spxextra{chowclassifier.chow.ChewData method}}

\begin{fulllineitems}
\phantomsection\label{\detokenize{source/chowclassifier:chowclassifier.chow.ChewData.parse_timecol}}
\pysigstartsignatures
\pysiglinewithargsret{\sphinxbfcode{\sphinxupquote{parse\_timecol}}}{\emph{\DUrole{n}{date\_format}\DUrole{p}{:}\DUrole{w}{  }\DUrole{n}{str}\DUrole{w}{  }\DUrole{o}{=}\DUrole{w}{  }\DUrole{default_value}{\textquotesingle{}\%Y\sphinxhyphen{}\%m\sphinxhyphen{}\%d\textquotesingle{}}}}{}
\pysigstopsignatures
\sphinxAtStartPar
Timecol must be numeric, if you have dates instead, use this method to transform
that column. (you need to adjust the breakpoint accordingly!)
\begin{quote}\begin{description}
\sphinxlineitem{Params format}
\sphinxAtStartPar
give the date format (default yyyy\sphinxhyphen{}mm\sphinxhyphen{}dd)

\end{description}\end{quote}

\end{fulllineitems}

\index{plot() (chowclassifier.chow.ChewData method)@\spxentry{plot()}\spxextra{chowclassifier.chow.ChewData method}}

\begin{fulllineitems}
\phantomsection\label{\detokenize{source/chowclassifier:chowclassifier.chow.ChewData.plot}}
\pysigstartsignatures
\pysiglinewithargsret{\sphinxbfcode{\sphinxupquote{plot}}}{\emph{\DUrole{n}{filename}\DUrole{p}{:}\DUrole{w}{  }\DUrole{n}{str}}, \emph{\DUrole{n}{title}\DUrole{p}{:}\DUrole{w}{  }\DUrole{n}{str}\DUrole{w}{  }\DUrole{o}{=}\DUrole{w}{  }\DUrole{default_value}{\textquotesingle{}\textquotesingle{}}}, \emph{\DUrole{n}{divide}\DUrole{p}{:}\DUrole{w}{  }\DUrole{n}{int}\DUrole{w}{  }\DUrole{o}{=}\DUrole{w}{  }\DUrole{default_value}{1}}, \emph{\DUrole{n}{nrows}\DUrole{p}{:}\DUrole{w}{  }\DUrole{n}{None\DUrole{w}{  }\DUrole{p}{|}\DUrole{w}{  }int}\DUrole{w}{  }\DUrole{o}{=}\DUrole{w}{  }\DUrole{default_value}{None}}, \emph{\DUrole{n}{ncols}\DUrole{p}{:}\DUrole{w}{  }\DUrole{n}{None\DUrole{w}{  }\DUrole{p}{|}\DUrole{w}{  }int}\DUrole{w}{  }\DUrole{o}{=}\DUrole{w}{  }\DUrole{default_value}{None}}, \emph{\DUrole{n}{scatterparams}\DUrole{p}{:}\DUrole{w}{  }\DUrole{n}{dict}\DUrole{w}{  }\DUrole{o}{=}\DUrole{w}{  }\DUrole{default_value}{\{\}}}, \emph{\DUrole{n}{fillparams}\DUrole{p}{:}\DUrole{w}{  }\DUrole{n}{dict}\DUrole{w}{  }\DUrole{o}{=}\DUrole{w}{  }\DUrole{default_value}{\{\}}}, \emph{\DUrole{n}{fitparams}\DUrole{p}{:}\DUrole{w}{  }\DUrole{n}{dict}\DUrole{w}{  }\DUrole{o}{=}\DUrole{w}{  }\DUrole{default_value}{\{\}}}, \emph{\DUrole{o}{**}\DUrole{n}{params}}}{}
\pysigstopsignatures
\sphinxAtStartPar
Plot all trends in one image
\begin{quote}\begin{description}
\sphinxlineitem{Params filename}
\sphinxAtStartPar
name of the output with extension (e.g name.pdf)

\sphinxlineitem{Params title}
\sphinxAtStartPar
title printed on the output

\sphinxlineitem{Params divide}
\sphinxAtStartPar
divide the plots into n files

\sphinxlineitem{Params ncols}
\sphinxAtStartPar
number of columns

\sphinxlineitem{Params nrows}
\sphinxAtStartPar
number of rows

\sphinxlineitem{Params scatterparams}
\sphinxAtStartPar
dictionnary of parameters for the scatter part of the plot, passed to chew plot

\sphinxlineitem{Params fillparams}
\sphinxAtStartPar
dictionnary of parameters for the error fill area of the plot, passed to chew plot

\sphinxlineitem{Params fitparams}
\sphinxAtStartPar
dictionnary of parameters for the fitted line, passed to chew plot

\sphinxlineitem{Params params}
\sphinxAtStartPar
named parameters passed to matplotlib, eg xlabel, ylabel, sharey, figsize, pad

\end{description}\end{quote}

\end{fulllineitems}

\index{plot\_by\_group() (chowclassifier.chow.ChewData method)@\spxentry{plot\_by\_group()}\spxextra{chowclassifier.chow.ChewData method}}

\begin{fulllineitems}
\phantomsection\label{\detokenize{source/chowclassifier:chowclassifier.chow.ChewData.plot_by_group}}
\pysigstartsignatures
\pysiglinewithargsret{\sphinxbfcode{\sphinxupquote{plot\_by\_group}}}{\emph{\DUrole{n}{groupcol}}, \emph{\DUrole{n}{savingpath}\DUrole{p}{:}\DUrole{w}{  }\DUrole{n}{str}\DUrole{w}{  }\DUrole{o}{=}\DUrole{w}{  }\DUrole{default_value}{\textquotesingle{}figs\textquotesingle{}}}, \emph{\DUrole{n}{format}\DUrole{p}{:}\DUrole{w}{  }\DUrole{n}{str}\DUrole{w}{  }\DUrole{o}{=}\DUrole{w}{  }\DUrole{default_value}{\textquotesingle{}pdf\textquotesingle{}}}, \emph{\DUrole{n}{cmap}\DUrole{p}{:}\DUrole{w}{  }\DUrole{n}{str}\DUrole{w}{  }\DUrole{o}{=}\DUrole{w}{  }\DUrole{default_value}{\textquotesingle{}Set1\textquotesingle{}}}, \emph{\DUrole{n}{show\_legend}\DUrole{p}{:}\DUrole{w}{  }\DUrole{n}{bool}\DUrole{w}{  }\DUrole{o}{=}\DUrole{w}{  }\DUrole{default_value}{True}}, \emph{\DUrole{o}{**}\DUrole{n}{params}}}{}
\pysigstopsignatures
\sphinxAtStartPar
Plot all trends in individual files
\begin{quote}\begin{description}
\sphinxlineitem{Params groupcol}
\sphinxAtStartPar
separate series by group, the analysis will be ran on each group seperately and toghether

\sphinxlineitem{Params savingpath}
\sphinxAtStartPar
path where the plots will be saved

\sphinxlineitem{Params format}
\sphinxAtStartPar
format of the image

\sphinxlineitem{Params cmap}
\sphinxAtStartPar
Name of one of matplotlib cmaps

\sphinxlineitem{Params show\_legend}
\sphinxAtStartPar
include the legend

\sphinxlineitem{Params params}
\sphinxAtStartPar
named parameters passed to matplotlib, eg. xlabel, ylabel, sharey, figsize, overall\_label

\end{description}\end{quote}

\end{fulllineitems}

\index{plot\_individually() (chowclassifier.chow.ChewData method)@\spxentry{plot\_individually()}\spxextra{chowclassifier.chow.ChewData method}}

\begin{fulllineitems}
\phantomsection\label{\detokenize{source/chowclassifier:chowclassifier.chow.ChewData.plot_individually}}
\pysigstartsignatures
\pysiglinewithargsret{\sphinxbfcode{\sphinxupquote{plot\_individually}}}{\emph{\DUrole{n}{savingpath}\DUrole{p}{:}\DUrole{w}{  }\DUrole{n}{str}\DUrole{w}{  }\DUrole{o}{=}\DUrole{w}{  }\DUrole{default_value}{\textquotesingle{}figs\textquotesingle{}}}, \emph{\DUrole{n}{format}\DUrole{p}{:}\DUrole{w}{  }\DUrole{n}{str}\DUrole{w}{  }\DUrole{o}{=}\DUrole{w}{  }\DUrole{default_value}{\textquotesingle{}pdf\textquotesingle{}}}, \emph{\DUrole{o}{**}\DUrole{n}{params}}}{}
\pysigstopsignatures
\sphinxAtStartPar
Plot all trends in individual files
\begin{quote}\begin{description}
\sphinxlineitem{Params savingpath}
\sphinxAtStartPar
path where the plots will be saved

\sphinxlineitem{Params format}
\sphinxAtStartPar
format of the image

\sphinxlineitem{Params params}
\sphinxAtStartPar
named parameters passed to matplotlib, eg. xlabel, ylabel, sharey, figsize

\end{description}\end{quote}

\end{fulllineitems}

\index{results\_df() (chowclassifier.chow.ChewData method)@\spxentry{results\_df()}\spxextra{chowclassifier.chow.ChewData method}}

\begin{fulllineitems}
\phantomsection\label{\detokenize{source/chowclassifier:chowclassifier.chow.ChewData.results_df}}
\pysigstartsignatures
\pysiglinewithargsret{\sphinxbfcode{\sphinxupquote{results\_df}}}{}{}
\pysigstopsignatures
\sphinxAtStartPar
Create dataframe with results

\end{fulllineitems}

\index{run() (chowclassifier.chow.ChewData method)@\spxentry{run()}\spxextra{chowclassifier.chow.ChewData method}}

\begin{fulllineitems}
\phantomsection\label{\detokenize{source/chowclassifier:chowclassifier.chow.ChewData.run}}
\pysigstartsignatures
\pysiglinewithargsret{\sphinxbfcode{\sphinxupquote{run}}}{\emph{\DUrole{n}{initial\_breakpoint}\DUrole{o}{=}\DUrole{default_value}{None}}}{}
\pysigstopsignatures
\sphinxAtStartPar
Run the classification for each group
\begin{quote}\begin{description}
\sphinxlineitem{Params initial\_breakpoint}
\sphinxAtStartPar
initial break point for Chow test (if None, will take midpoint of timecol)

\end{description}\end{quote}

\end{fulllineitems}

\index{run\_C() (chowclassifier.chow.ChewData method)@\spxentry{run\_C()}\spxextra{chowclassifier.chow.ChewData method}}

\begin{fulllineitems}
\phantomsection\label{\detokenize{source/chowclassifier:chowclassifier.chow.ChewData.run_C}}
\pysigstartsignatures
\pysiglinewithargsret{\sphinxbfcode{\sphinxupquote{run\_C}}}{\emph{\DUrole{n}{sub\_df: \textasciitilde{}pandas.core.frame.DataFrame}}, \emph{\DUrole{n}{name: str}}, \emph{\DUrole{n}{initial\_breakpoint: float | int) \sphinxhyphen{}\textgreater{} (\textless{}class \textquotesingle{}dict\textquotesingle{}\textgreater{}}}, \emph{\DUrole{n}{\textless{}class \textquotesingle{}chowclassifier.chow.Chow\textquotesingle{}\textgreater{}}}}{}
\pysigstopsignatures
\sphinxAtStartPar
Run the classification on a dataframe
\begin{quote}\begin{description}
\sphinxlineitem{Params sub\_df}
\sphinxAtStartPar
dataframe

\sphinxlineitem{Params name}
\sphinxAtStartPar
name of the series

\sphinxlineitem{Params initial\_breakpoint}
\sphinxAtStartPar
initial break point for Chow test (if None, will take midpoint of timecol)

\sphinxlineitem{Returns}
\sphinxAtStartPar
dictionary of parameters from the analysis, Chow object from the analysis

\end{description}\end{quote}

\end{fulllineitems}

\index{save\_to\_file() (chowclassifier.chow.ChewData method)@\spxentry{save\_to\_file()}\spxextra{chowclassifier.chow.ChewData method}}

\begin{fulllineitems}
\phantomsection\label{\detokenize{source/chowclassifier:chowclassifier.chow.ChewData.save_to_file}}
\pysigstartsignatures
\pysiglinewithargsret{\sphinxbfcode{\sphinxupquote{save\_to\_file}}}{\emph{\DUrole{n}{filename}\DUrole{p}{:}\DUrole{w}{  }\DUrole{n}{str}}, \emph{\DUrole{o}{**}\DUrole{n}{kwargs}}}{}
\pysigstopsignatures
\sphinxAtStartPar
Export results to excel or csv
\begin{quote}\begin{description}
\sphinxlineitem{Params filename}
\sphinxAtStartPar
filename

\sphinxlineitem{Params kwargs}
\sphinxAtStartPar
other parameters passed to pd.DataFrame.to\_csv

\end{description}\end{quote}

\end{fulllineitems}

\index{to\_csv() (chowclassifier.chow.ChewData method)@\spxentry{to\_csv()}\spxextra{chowclassifier.chow.ChewData method}}

\begin{fulllineitems}
\phantomsection\label{\detokenize{source/chowclassifier:chowclassifier.chow.ChewData.to_csv}}
\pysigstartsignatures
\pysiglinewithargsret{\sphinxbfcode{\sphinxupquote{to\_csv}}}{\emph{\DUrole{n}{filename}\DUrole{p}{:}\DUrole{w}{  }\DUrole{n}{str}}, \emph{\DUrole{o}{**}\DUrole{n}{kwargs}}}{}
\pysigstopsignatures
\sphinxAtStartPar
Export results to csv
\begin{quote}\begin{description}
\sphinxlineitem{Params filename}
\sphinxAtStartPar
filename

\end{description}\end{quote}

\end{fulllineitems}

\index{to\_excel() (chowclassifier.chow.ChewData method)@\spxentry{to\_excel()}\spxextra{chowclassifier.chow.ChewData method}}

\begin{fulllineitems}
\phantomsection\label{\detokenize{source/chowclassifier:chowclassifier.chow.ChewData.to_excel}}
\pysigstartsignatures
\pysiglinewithargsret{\sphinxbfcode{\sphinxupquote{to\_excel}}}{\emph{\DUrole{n}{filename}\DUrole{p}{:}\DUrole{w}{  }\DUrole{n}{str}}, \emph{\DUrole{o}{**}\DUrole{n}{kwargs}}}{}
\pysigstopsignatures
\sphinxAtStartPar
Export results to excel
\begin{quote}\begin{description}
\sphinxlineitem{Params filename}
\sphinxAtStartPar
filename

\end{description}\end{quote}

\end{fulllineitems}


\end{fulllineitems}

\index{Chow (class in chowclassifier.chow)@\spxentry{Chow}\spxextra{class in chowclassifier.chow}}

\begin{fulllineitems}
\phantomsection\label{\detokenize{source/chowclassifier:chowclassifier.chow.Chow}}
\pysigstartsignatures
\pysiglinewithargsret{\sphinxbfcode{\sphinxupquote{class\DUrole{w}{  }}}\sphinxcode{\sphinxupquote{chowclassifier.chow.}}\sphinxbfcode{\sphinxupquote{Chow}}}{\emph{\DUrole{n}{df}}, \emph{\DUrole{n}{name}\DUrole{p}{:}\DUrole{w}{  }\DUrole{n}{str}\DUrole{w}{  }\DUrole{o}{=}\DUrole{w}{  }\DUrole{default_value}{\textquotesingle{}\textquotesingle{}}}, \emph{\DUrole{n}{initial\_breakpoint}\DUrole{p}{:}\DUrole{w}{  }\DUrole{n}{float\DUrole{w}{  }\DUrole{p}{|}\DUrole{w}{  }None}\DUrole{w}{  }\DUrole{o}{=}\DUrole{w}{  }\DUrole{default_value}{None}}, \emph{\DUrole{n}{breakpoints}\DUrole{p}{:}\DUrole{w}{  }\DUrole{n}{list\DUrole{w}{  }\DUrole{p}{|}\DUrole{w}{  }None}\DUrole{w}{  }\DUrole{o}{=}\DUrole{w}{  }\DUrole{default_value}{None}}, \emph{\DUrole{n}{timecol}\DUrole{p}{:}\DUrole{w}{  }\DUrole{n}{str}\DUrole{w}{  }\DUrole{o}{=}\DUrole{w}{  }\DUrole{default_value}{\textquotesingle{}year\textquotesingle{}}}, \emph{\DUrole{n}{ycol}\DUrole{p}{:}\DUrole{w}{  }\DUrole{n}{str}\DUrole{w}{  }\DUrole{o}{=}\DUrole{w}{  }\DUrole{default_value}{\textquotesingle{}value\textquotesingle{}}}, \emph{\DUrole{n}{alpha}\DUrole{o}{=}\DUrole{default_value}{0.01}}, \emph{\DUrole{n}{margin}\DUrole{p}{:}\DUrole{w}{  }\DUrole{n}{int}\DUrole{w}{  }\DUrole{o}{=}\DUrole{w}{  }\DUrole{default_value}{2}}}{}
\pysigstopsignatures
\sphinxAtStartPar
Bases: \sphinxcode{\sphinxupquote{object}}

\sphinxAtStartPar
Class handling the chow analysis
Call \sphinxtitleref{run()} to run the analysis (find breakpoint and classify)
Call afterwards \sphinxtitleref{params()} to get the regressions (y=k * x + b) parameters as
{[}{[}model1\_k, model1\_b, model1\_Rsquared{]},{[}model2\_k, model2\_b, model2\_Rsquared{]}{]}
(if the breakpoint is not significant, model1 = model2)
The significance level of the breakpoint is adjusted with Bonferroni correction,
i.e. dividing alpha by the number of tests. The confidence interval of the slope
and intercept are at significance level alpha (not corrected).
\begin{quote}\begin{description}
\sphinxlineitem{Params df}
\sphinxAtStartPar
data on which to perform the Chow categorisation

\sphinxlineitem{Params name}
\sphinxAtStartPar
name of the dataset (optional)

\sphinxlineitem{Params breakpoint}
\sphinxAtStartPar
breakpoint (if None, midpoint of X series will be used)

\sphinxlineitem{Params breakpoints}
\sphinxAtStartPar
candidates for the breakpoints

\sphinxlineitem{Params timecol}
\sphinxAtStartPar
name of time column (X)

\sphinxlineitem{Params ycol}
\sphinxAtStartPar
name of the variable column (Y)

\sphinxlineitem{Parameters}
\sphinxAtStartPar
\sphinxstyleliteralstrong{\sphinxupquote{alpha}} \textendash{} significance level (without Bonferroni correction), can be changed with \sphinxcode{\sphinxupquote{Chow.set\_alpha(0.01)}}

\sphinxlineitem{Params margin}
\sphinxAtStartPar
range of indexes around breakpoint (or midpoint) where the best breakpoint
will be searched. E.g. if X = {[}0,1,1.5,2,3,4,5{]}, breakpoint = 2 and margin = 1,
the set of breakpoints {[}1.5,2,3{]} will be tested.
(set = 0 if you want to look only at the given breakpoint or the midpoint)

\end{description}\end{quote}
\index{classify() (chowclassifier.chow.Chow method)@\spxentry{classify()}\spxextra{chowclassifier.chow.Chow method}}

\begin{fulllineitems}
\phantomsection\label{\detokenize{source/chowclassifier:chowclassifier.chow.Chow.classify}}
\pysigstartsignatures
\pysiglinewithargsret{\sphinxbfcode{\sphinxupquote{classify}}}{\emph{\DUrole{o}{**}\DUrole{n}{kwargs}}}{}
\pysigstopsignatures
\sphinxAtStartPar
classify in the dataset based on results of Chow test at significance level alpha
Possible values are:
No significant breakpoint and:
\begin{enumerate}
\sphinxsetlistlabels{\arabic}{enumi}{enumii}{}{.}%
\item {} 
\sphinxAtStartPar
N: non\sphinxhyphen{}significant overall trend

\item {} 
\sphinxAtStartPar
I: significant increasing overall trend

\item {} 
\sphinxAtStartPar
D: significant decreasing overall trend

\end{enumerate}

\sphinxAtStartPar
Significant breakpoint (set1 and set2 indicate points before, respectively after, breakpoint) and:
\begin{enumerate}
\sphinxsetlistlabels{\arabic}{enumi}{enumii}{}{.}%
\setcounter{enumi}{3}
\item {} 
\sphinxAtStartPar
NN: non\sphinxhyphen{}significant trend on set1 and non\sphinxhyphen{}significant trend on set2

\item {} 
\sphinxAtStartPar
NI: non\sphinxhyphen{}significant trend on set1 and significant increasing trend on set2

\item {} 
\sphinxAtStartPar
ND: non\sphinxhyphen{}significant trend on set1 ans significant decreasing trend on set2

\item {} 
\sphinxAtStartPar
IN: significant increasing trend on set1 and non\sphinxhyphen{}significant trend on set2

\item {} 
\sphinxAtStartPar
ID: significant increasing trend on set1 and significant decreasing trend on set2

\item {} 
\sphinxAtStartPar
iI: significant increasing trend on both set1 and set2 with greater increase in set2

\item {} 
\sphinxAtStartPar
Ii: significant increasing trend on both set1 and set2 with greater increase in set1

\item {} 
\sphinxAtStartPar
DN: significant decreasing trend on set1 and non\sphinxhyphen{}significant trend on set2

\item {} 
\sphinxAtStartPar
DI: significant decreasing trend on set1 and significant increasing trend on set2

\item {} 
\sphinxAtStartPar
dD: significant decreasing trend on both set1 and set2 with greater decrease in set2

\item {} 
\sphinxAtStartPar
Dd: significant decreasing trend on both set1 and set2 with greater decrease in set1

\end{enumerate}

\end{fulllineitems}

\index{find\_best\_bkp() (chowclassifier.chow.Chow method)@\spxentry{find\_best\_bkp()}\spxextra{chowclassifier.chow.Chow method}}

\begin{fulllineitems}
\phantomsection\label{\detokenize{source/chowclassifier:chowclassifier.chow.Chow.find_best_bkp}}
\pysigstartsignatures
\pysiglinewithargsret{\sphinxbfcode{\sphinxupquote{find\_best\_bkp}}}{\emph{\DUrole{n}{breakpoint\_indices}\DUrole{o}{=}\DUrole{default_value}{None}}, \emph{\DUrole{o}{**}\DUrole{n}{kwargs}}}{}
\pysigstopsignatures
\sphinxAtStartPar
finds best breakpoint and run Chow test and find OLS params for either whole set,
and the two subset broken at best breakpoint

\end{fulllineitems}

\index{get\_breakpoint\_indices() (chowclassifier.chow.Chow method)@\spxentry{get\_breakpoint\_indices()}\spxextra{chowclassifier.chow.Chow method}}

\begin{fulllineitems}
\phantomsection\label{\detokenize{source/chowclassifier:chowclassifier.chow.Chow.get_breakpoint_indices}}
\pysigstartsignatures
\pysiglinewithargsret{\sphinxbfcode{\sphinxupquote{get\_breakpoint\_indices}}}{}{}
\pysigstopsignatures
\sphinxAtStartPar
sets the breakpoints

\end{fulllineitems}

\index{params() (chowclassifier.chow.Chow method)@\spxentry{params()}\spxextra{chowclassifier.chow.Chow method}}

\begin{fulllineitems}
\phantomsection\label{\detokenize{source/chowclassifier:chowclassifier.chow.Chow.params}}
\pysigstartsignatures
\pysiglinewithargsret{\sphinxbfcode{\sphinxupquote{params}}}{}{}
\pysigstopsignatures
\sphinxAtStartPar
Returns parameters of model

\end{fulllineitems}

\index{params\_names() (chowclassifier.chow.Chow method)@\spxentry{params\_names()}\spxextra{chowclassifier.chow.Chow method}}

\begin{fulllineitems}
\phantomsection\label{\detokenize{source/chowclassifier:chowclassifier.chow.Chow.params_names}}
\pysigstartsignatures
\pysiglinewithargsret{\sphinxbfcode{\sphinxupquote{params\_names}}}{}{}
\pysigstopsignatures
\end{fulllineitems}

\index{plot() (chowclassifier.chow.Chow method)@\spxentry{plot()}\spxextra{chowclassifier.chow.Chow method}}

\begin{fulllineitems}
\phantomsection\label{\detokenize{source/chowclassifier:chowclassifier.chow.Chow.plot}}
\pysigstartsignatures
\pysiglinewithargsret{\sphinxbfcode{\sphinxupquote{plot}}}{\emph{\DUrole{n}{color}\DUrole{o}{=}\DUrole{default_value}{None}}, \emph{\DUrole{n}{filename}\DUrole{p}{:}\DUrole{w}{  }\DUrole{n}{str\DUrole{w}{  }\DUrole{p}{|}\DUrole{w}{  }None}\DUrole{w}{  }\DUrole{o}{=}\DUrole{w}{  }\DUrole{default_value}{None}}, \emph{\DUrole{n}{ax}\DUrole{o}{=}\DUrole{default_value}{None}}, \emph{\DUrole{n}{figsize}\DUrole{o}{=}\DUrole{default_value}{(16, 8)}}, \emph{\DUrole{n}{ylabel}\DUrole{p}{:}\DUrole{w}{  }\DUrole{n}{str}\DUrole{w}{  }\DUrole{o}{=}\DUrole{w}{  }\DUrole{default_value}{\textquotesingle{}\textquotesingle{}}}, \emph{\DUrole{n}{xlabel}\DUrole{p}{:}\DUrole{w}{  }\DUrole{n}{str}\DUrole{w}{  }\DUrole{o}{=}\DUrole{w}{  }\DUrole{default_value}{\textquotesingle{}\textquotesingle{}}}, \emph{\DUrole{n}{title}\DUrole{p}{:}\DUrole{w}{  }\DUrole{n}{str}\DUrole{w}{  }\DUrole{o}{=}\DUrole{w}{  }\DUrole{default_value}{\textquotesingle{}\textquotesingle{}}}, \emph{\DUrole{n}{ylog}\DUrole{p}{:}\DUrole{w}{  }\DUrole{n}{bool}\DUrole{w}{  }\DUrole{o}{=}\DUrole{w}{  }\DUrole{default_value}{False}}, \emph{\DUrole{n}{fill}\DUrole{p}{:}\DUrole{w}{  }\DUrole{n}{bool}\DUrole{w}{  }\DUrole{o}{=}\DUrole{w}{  }\DUrole{default_value}{True}}, \emph{\DUrole{n}{scatterparams}\DUrole{p}{:}\DUrole{w}{  }\DUrole{n}{dict}\DUrole{w}{  }\DUrole{o}{=}\DUrole{w}{  }\DUrole{default_value}{\{\}}}, \emph{\DUrole{n}{fillparams}\DUrole{p}{:}\DUrole{w}{  }\DUrole{n}{dict}\DUrole{w}{  }\DUrole{o}{=}\DUrole{w}{  }\DUrole{default_value}{\{\}}}, \emph{\DUrole{n}{fitparams}\DUrole{p}{:}\DUrole{w}{  }\DUrole{n}{dict}\DUrole{w}{  }\DUrole{o}{=}\DUrole{w}{  }\DUrole{default_value}{\{\}}}}{}
\pysigstopsignatures
\sphinxAtStartPar
Plot the regression model(s)
\begin{quote}\begin{description}
\sphinxlineitem{Params color}
\sphinxAtStartPar
give color for the points (default defined in TREND\_COLORS)

\sphinxlineitem{Params filename}
\sphinxAtStartPar
filename (with extension format), if None, will call plt.show() instead

\sphinxlineitem{Params ax}
\sphinxAtStartPar
ax on which to plot, if None will create new figure

\sphinxlineitem{Params figsize}
\sphinxAtStartPar
tuple with figure size

\sphinxlineitem{Params ylabel}
\sphinxAtStartPar
label for y\sphinxhyphen{}axis

\sphinxlineitem{Params xlabel}
\sphinxAtStartPar
label for x\sphinxhyphen{}axis

\sphinxlineitem{Params title}
\sphinxAtStartPar
title of the plot/subplot

\sphinxlineitem{Params ylog}
\sphinxAtStartPar
make the y\sphinxhyphen{}axis log\sphinxhyphen{}scale

\sphinxlineitem{Params fill}
\sphinxAtStartPar
print the confidence interval

\sphinxlineitem{Params scatterparams}
\sphinxAtStartPar
parameters for the scatter plot

\sphinxlineitem{Params fillparams}
\sphinxAtStartPar
parameter for the confidence intervals

\sphinxlineitem{Params fitparams}
\sphinxAtStartPar
parameter for the fitted line

\end{description}\end{quote}

\end{fulllineitems}

\index{run() (chowclassifier.chow.Chow method)@\spxentry{run()}\spxextra{chowclassifier.chow.Chow method}}

\begin{fulllineitems}
\phantomsection\label{\detokenize{source/chowclassifier:chowclassifier.chow.Chow.run}}
\pysigstartsignatures
\pysiglinewithargsret{\sphinxbfcode{\sphinxupquote{run}}}{\emph{\DUrole{o}{**}\DUrole{n}{kwargs}}}{}
\pysigstopsignatures
\sphinxAtStartPar
Run the chow classification. parameters will be forwarded to run and classify,
e.g. initial\_breakpoint can be passed to run, and alpha to classify

\end{fulllineitems}

\index{run\_chow() (chowclassifier.chow.Chow method)@\spxentry{run\_chow()}\spxextra{chowclassifier.chow.Chow method}}

\begin{fulllineitems}
\phantomsection\label{\detokenize{source/chowclassifier:chowclassifier.chow.Chow.run_chow}}
\pysigstartsignatures
\pysiglinewithargsret{\sphinxbfcode{\sphinxupquote{run\_chow}}}{\emph{\DUrole{n}{bkp}}, \emph{\DUrole{n}{normalise}\DUrole{o}{=}\DUrole{default_value}{False}}}{}
\pysigstopsignatures
\sphinxAtStartPar
Run chow test at breakpoint

\end{fulllineitems}

\index{set\_alpha() (chowclassifier.chow.Chow method)@\spxentry{set\_alpha()}\spxextra{chowclassifier.chow.Chow method}}

\begin{fulllineitems}
\phantomsection\label{\detokenize{source/chowclassifier:chowclassifier.chow.Chow.set_alpha}}
\pysigstartsignatures
\pysiglinewithargsret{\sphinxbfcode{\sphinxupquote{set\_alpha}}}{\emph{\DUrole{n}{alpha}\DUrole{p}{:}\DUrole{w}{  }\DUrole{n}{float}}}{}
\pysigstopsignatures
\sphinxAtStartPar
Changes the current statistical significance level
The bonferroni correction will be applied
\begin{quote}\begin{description}
\sphinxlineitem{Params alpha}
\sphinxAtStartPar
positive float in {]}0,1{[}

\end{description}\end{quote}

\end{fulllineitems}

\index{set\_color() (chowclassifier.chow.Chow method)@\spxentry{set\_color()}\spxextra{chowclassifier.chow.Chow method}}

\begin{fulllineitems}
\phantomsection\label{\detokenize{source/chowclassifier:chowclassifier.chow.Chow.set_color}}
\pysigstartsignatures
\pysiglinewithargsret{\sphinxbfcode{\sphinxupquote{set\_color}}}{}{}
\pysigstopsignatures
\end{fulllineitems}

\index{summary() (chowclassifier.chow.Chow method)@\spxentry{summary()}\spxextra{chowclassifier.chow.Chow method}}

\begin{fulllineitems}
\phantomsection\label{\detokenize{source/chowclassifier:chowclassifier.chow.Chow.summary}}
\pysigstartsignatures
\pysiglinewithargsret{\sphinxbfcode{\sphinxupquote{summary}}}{}{}
\pysigstopsignatures
\sphinxAtStartPar
Print the summary of the model(s)

\end{fulllineitems}


\end{fulllineitems}

\index{conf\_int\_to\_trend() (in module chowclassifier.chow)@\spxentry{conf\_int\_to\_trend()}\spxextra{in module chowclassifier.chow}}

\begin{fulllineitems}
\phantomsection\label{\detokenize{source/chowclassifier:chowclassifier.chow.conf_int_to_trend}}
\pysigstartsignatures
\pysiglinewithargsret{\sphinxcode{\sphinxupquote{chowclassifier.chow.}}\sphinxbfcode{\sphinxupquote{conf\_int\_to\_trend}}}{\emph{\DUrole{n}{ci}\DUrole{p}{:}\DUrole{w}{  }\DUrole{n}{DataFrame\DUrole{w}{  }\DUrole{p}{|}\DUrole{w}{  }list}}, \emph{\DUrole{n}{param}\DUrole{p}{:}\DUrole{w}{  }\DUrole{n}{str}\DUrole{w}{  }\DUrole{o}{=}\DUrole{w}{  }\DUrole{default_value}{\textquotesingle{}X\textquotesingle{}}}}{}
\pysigstopsignatures
\sphinxAtStartPar
return trend  0, 1, 2 from the confidence interval:
\begin{itemize}
\item {} 
\sphinxAtStartPar
if ci crosses 0, then returns 0

\item {} 
\sphinxAtStartPar
if ci is completely below 0, then returns 2

\item {} 
\sphinxAtStartPar
else returns 1

\end{itemize}
\begin{quote}\begin{description}
\sphinxlineitem{Params ci}
\sphinxAtStartPar
confidence interval to process, either a tuple (length 2), a dataframe or result from conf\_int given by statsmodles

\sphinxlineitem{Params param}
\sphinxAtStartPar
name column to use if a dataframe was given as ci

\end{description}\end{quote}

\end{fulllineitems}

\index{guess\_timelapse\_type() (in module chowclassifier.chow)@\spxentry{guess\_timelapse\_type()}\spxextra{in module chowclassifier.chow}}

\begin{fulllineitems}
\phantomsection\label{\detokenize{source/chowclassifier:chowclassifier.chow.guess_timelapse_type}}
\pysigstartsignatures
\pysiglinewithargsret{\sphinxcode{\sphinxupquote{chowclassifier.chow.}}\sphinxbfcode{\sphinxupquote{guess\_timelapse\_type}}}{\emph{\DUrole{n}{series}\DUrole{p}{:}\DUrole{w}{  }\DUrole{n}{list\DUrole{w}{  }\DUrole{p}{|}\DUrole{w}{  }DataFrame\DUrole{w}{  }\DUrole{p}{|}\DUrole{w}{  }Series}}, \emph{\DUrole{n}{date\_format}\DUrole{p}{:}\DUrole{w}{  }\DUrole{n}{str}\DUrole{w}{  }\DUrole{o}{=}\DUrole{w}{  }\DUrole{default_value}{\textquotesingle{}\%Y\sphinxhyphen{}\%m\sphinxhyphen{}\%d\textquotesingle{}}}}{}
\pysigstopsignatures
\sphinxAtStartPar
Guess the time lapse in a series and index it
\begin{quote}\begin{description}
\sphinxlineitem{Params series}
\sphinxAtStartPar
data

\sphinxlineitem{Params date\_format}
\sphinxAtStartPar
format of the datetime

\sphinxlineitem{Returns}
\sphinxAtStartPar
indexed series and datetype

\end{description}\end{quote}

\end{fulllineitems}

\index{infer\_date\_format() (in module chowclassifier.chow)@\spxentry{infer\_date\_format()}\spxextra{in module chowclassifier.chow}}

\begin{fulllineitems}
\phantomsection\label{\detokenize{source/chowclassifier:chowclassifier.chow.infer_date_format}}
\pysigstartsignatures
\pysiglinewithargsret{\sphinxcode{\sphinxupquote{chowclassifier.chow.}}\sphinxbfcode{\sphinxupquote{infer\_date\_format}}}{\emph{\DUrole{n}{s\_date}\DUrole{p}{:}\DUrole{w}{  }\DUrole{n}{str}}}{}
\pysigstopsignatures
\sphinxAtStartPar
Infer date format from a sample
\begin{quote}\begin{description}
\sphinxlineitem{Params s\_date}
\sphinxAtStartPar
sample date to infer date format from.

\end{description}\end{quote}

\end{fulllineitems}



\subsubsection{Module contents}
\label{\detokenize{source/chowclassifier:module-chowclassifier}}\label{\detokenize{source/chowclassifier:module-contents}}\index{module@\spxentry{module}!chowclassifier@\spxentry{chowclassifier}}\index{chowclassifier@\spxentry{chowclassifier}!module@\spxentry{module}}

\chapter{Licence}
\label{\detokenize{index:licence}}
\sphinxAtStartPar
GNU LESSER GENERAL PUBLIC LICENSE
Version 2.1, February 1999


\chapter{Outputs}
\label{\detokenize{index:outputs}}

\section{\sphinxstyleliteralintitle{\sphinxupquote{plot}} and \sphinxstyleliteralintitle{\sphinxupquote{plot\_individually}}}
\label{\detokenize{index:plot-and-plot-individually}}
\sphinxAtStartPar
Both methods, plot the best models (either the regression on the whole dataset or, two separate regressions on the left and right parts of the dataset split at breakpoint) for all subdatasets (either each in individual files or together).
The filled area represents the confidence interval for the predicted values in the regression model/models. The scatter points show the observed data. The lines are the regression lines.


\section{\sphinxstyleliteralintitle{\sphinxupquote{save\_to\_file}}}
\label{\detokenize{index:save-to-file}}
\sphinxAtStartPar
The output’s columns are described in the following table. \sphinxcode{\sphinxupquote{model0}} correspond to the regression on the whole dataset. \sphinxcode{\sphinxupquote{model1}} and \sphinxcode{\sphinxupquote{model2}} correspond to regressions on respectively the left and right parts of the dataset split at the breakpoint.


\begin{savenotes}
\sphinxatlongtablestart
\sphinxthistablewithglobalstyle
\begin{longtable}[c]{lll}
\sphinxtoprule
\sphinxstyletheadfamily 
\sphinxAtStartPar
Model
&\sphinxstyletheadfamily 
\sphinxAtStartPar
property
&\sphinxstyletheadfamily 
\sphinxAtStartPar
Description of the property
\\
\sphinxmidrule
\endfirsthead

\multicolumn{3}{c}{\sphinxnorowcolor
    \makebox[0pt]{\sphinxtablecontinued{\tablename\ \thetable{} \textendash{} continued from previous page}}%
}\\
\sphinxtoprule
\sphinxstyletheadfamily 
\sphinxAtStartPar
Model
&\sphinxstyletheadfamily 
\sphinxAtStartPar
property
&\sphinxstyletheadfamily 
\sphinxAtStartPar
Description of the property
\\
\sphinxmidrule
\endhead

\sphinxbottomrule
\multicolumn{3}{r}{\sphinxnorowcolor
    \makebox[0pt][r]{\sphinxtablecontinued{continues on next page}}%
}\\
\endfoot

\endlastfoot
\sphinxtableatstartofbodyhook

\sphinxAtStartPar
Global
&
\sphinxAtStartPar
name
&
\sphinxAtStartPar
Name of the group
\\
\sphinxhline
\sphinxAtStartPar
Global
&
\sphinxAtStartPar
breakpoint
&
\sphinxAtStartPar
Significant breakpoint (if found)
\\
\sphinxhline
\sphinxAtStartPar
Global
&
\sphinxAtStartPar
score
&
\sphinxAtStartPar
lowest p\sphinxhyphen{}value over the tested breakpoints
\\
\sphinxhline
\sphinxAtStartPar
model0
&
\sphinxAtStartPar
intercept
&
\sphinxAtStartPar
intercept
\\
\sphinxhline
\sphinxAtStartPar
model0
&
\sphinxAtStartPar
slope
&
\sphinxAtStartPar
slope
\\
\sphinxhline
\sphinxAtStartPar
model0
&
\sphinxAtStartPar
R2
&
\sphinxAtStartPar
coefficient of determination
\\
\sphinxhline
\sphinxAtStartPar
model0
&
\sphinxAtStartPar
se\_intercept
&
\sphinxAtStartPar
standard error for the intercept
\\
\sphinxhline
\sphinxAtStartPar
model0
&
\sphinxAtStartPar
se\_slope
&
\sphinxAtStartPar
standard error for the slope
\\
\sphinxhline
\sphinxAtStartPar
model0
&
\sphinxAtStartPar
n
&
\sphinxAtStartPar
number of data points used in the fit
\\
\sphinxhline
\sphinxAtStartPar
model0
&
\sphinxAtStartPar
intercept\_pvalue
&
\sphinxAtStartPar
p\sphinxhyphen{}value of the intercept
\\
\sphinxhline
\sphinxAtStartPar
model0
&
\sphinxAtStartPar
slope\_pvalue
&
\sphinxAtStartPar
p\sphinxhyphen{}value of the slope
\\
\sphinxhline
\sphinxAtStartPar
model0
&
\sphinxAtStartPar
intercept\_ci\_lower
&
\sphinxAtStartPar
lower bound of the confidence interval of the intercept
\\
\sphinxhline
\sphinxAtStartPar
model0
&
\sphinxAtStartPar
intercept\_ci\_upper
&
\sphinxAtStartPar
upper bound of the confidence interval of the intercept
\\
\sphinxhline
\sphinxAtStartPar
model0
&
\sphinxAtStartPar
slope\_ci\_lower
&
\sphinxAtStartPar
lower bound of the confidence slope of the intercept
\\
\sphinxhline
\sphinxAtStartPar
model0
&
\sphinxAtStartPar
slope\_ci\_upper
&
\sphinxAtStartPar
upper bound of the confidence interval of the slope
\\
\sphinxhline&
\sphinxAtStartPar
alpha
&
\sphinxAtStartPar
significance level at which the regressions were tested
\\
\sphinxhline&
\sphinxAtStartPar
alpha\_Bonferroni\_corrected
&
\sphinxAtStartPar
Bonferroni corrected significance level at which the breakpoints were tested
\\
\sphinxhline
\sphinxAtStartPar
model1
&
\sphinxAtStartPar
intercept
&
\sphinxAtStartPar
intercept
\\
\sphinxhline
\sphinxAtStartPar
model1
&
\sphinxAtStartPar
slope
&
\sphinxAtStartPar
slope
\\
\sphinxhline
\sphinxAtStartPar
model1
&
\sphinxAtStartPar
R2
&
\sphinxAtStartPar
coefficient of determination
\\
\sphinxhline
\sphinxAtStartPar
model1
&
\sphinxAtStartPar
se\_intercept
&
\sphinxAtStartPar
standard error for the intercept
\\
\sphinxhline
\sphinxAtStartPar
model1
&
\sphinxAtStartPar
se\_slope
&
\sphinxAtStartPar
standard error for the slope
\\
\sphinxhline
\sphinxAtStartPar
model1
&
\sphinxAtStartPar
n
&
\sphinxAtStartPar
number of data points used in the fit
\\
\sphinxhline
\sphinxAtStartPar
model1
&
\sphinxAtStartPar
intercept\_pvalue
&
\sphinxAtStartPar
p\sphinxhyphen{}value of the intercept
\\
\sphinxhline
\sphinxAtStartPar
model1
&
\sphinxAtStartPar
slope\_pvalue
&
\sphinxAtStartPar
p\sphinxhyphen{}value of the slope
\\
\sphinxhline
\sphinxAtStartPar
model1
&
\sphinxAtStartPar
intercept\_ci\_lower
&
\sphinxAtStartPar
lower bound of the confidence interval of the intercept
\\
\sphinxhline
\sphinxAtStartPar
model1
&
\sphinxAtStartPar
intercept\_ci\_upper
&
\sphinxAtStartPar
upper bound of the confidence interval of the intercept
\\
\sphinxhline
\sphinxAtStartPar
model1
&
\sphinxAtStartPar
slope\_ci\_lower
&
\sphinxAtStartPar
lower bound of the confidence slope of the intercept
\\
\sphinxhline
\sphinxAtStartPar
model1
&
\sphinxAtStartPar
slope\_ci\_upper
&
\sphinxAtStartPar
upper bound of the confidence interval of the slope
\\
\sphinxhline&
\sphinxAtStartPar
alpha
&
\sphinxAtStartPar
significance level at which the regressions were tested
\\
\sphinxhline&
\sphinxAtStartPar
alpha\_Bonferroni\_corrected
&
\sphinxAtStartPar
Bonferroni corrected significance level at which the breakpoints were tested
\\
\sphinxhline
\sphinxAtStartPar
model2
&
\sphinxAtStartPar
intercept
&
\sphinxAtStartPar
intercept
\\
\sphinxhline
\sphinxAtStartPar
model2
&
\sphinxAtStartPar
slope
&
\sphinxAtStartPar
slope
\\
\sphinxhline
\sphinxAtStartPar
model2
&
\sphinxAtStartPar
R2
&
\sphinxAtStartPar
coefficient of determination
\\
\sphinxhline
\sphinxAtStartPar
model2
&
\sphinxAtStartPar
se\_intercept
&
\sphinxAtStartPar
standard error for the intercept
\\
\sphinxhline
\sphinxAtStartPar
model2
&
\sphinxAtStartPar
se\_slope
&
\sphinxAtStartPar
standard error for the slope
\\
\sphinxhline
\sphinxAtStartPar
model2
&
\sphinxAtStartPar
n
&
\sphinxAtStartPar
number of data points used in the fit
\\
\sphinxhline
\sphinxAtStartPar
model2
&
\sphinxAtStartPar
intercept\_pvalue
&
\sphinxAtStartPar
p\sphinxhyphen{}value of the intercept
\\
\sphinxhline
\sphinxAtStartPar
model2
&
\sphinxAtStartPar
slope\_pvalue
&
\sphinxAtStartPar
p\sphinxhyphen{}value of the slope
\\
\sphinxhline
\sphinxAtStartPar
model2
&
\sphinxAtStartPar
intercept\_ci\_lower
&
\sphinxAtStartPar
lower bound of the confidence interval of the intercept
\\
\sphinxhline
\sphinxAtStartPar
model2
&
\sphinxAtStartPar
intercept\_ci\_upper
&
\sphinxAtStartPar
upper bound of the confidence interval of the intercept
\\
\sphinxhline
\sphinxAtStartPar
model2
&
\sphinxAtStartPar
slope\_ci\_lower
&
\sphinxAtStartPar
lower bound of the confidence slope of the intercept
\\
\sphinxhline
\sphinxAtStartPar
model2
&
\sphinxAtStartPar
slope\_ci\_upper
&
\sphinxAtStartPar
upper bound of the confidence interval of the slope
\\
\sphinxhline&
\sphinxAtStartPar
alpha
&
\sphinxAtStartPar
significance level at which the regressions were tested
\\
\sphinxhline&
\sphinxAtStartPar
alpha\_Bonferroni\_corrected
&
\sphinxAtStartPar
Bonferroni corrected significance level at which the breakpoints were tested
\\
\sphinxbottomrule
\end{longtable}
\sphinxtableafterendhook
\sphinxatlongtableend
\end{savenotes}


\chapter{Version changes}
\label{\detokenize{index:version-changes}}

\begin{savenotes}\sphinxattablestart
\sphinxthistablewithglobalstyle
\centering
\begin{tabulary}{\linewidth}[t]{TT}
\sphinxtoprule
\sphinxtableatstartofbodyhook
\sphinxAtStartPar
1.0.5
&
\sphinxAtStartPar
Added option to plot confidence interval fill for individual group trends
\\
\sphinxhline
\sphinxAtStartPar
1.0.4
&
\sphinxAtStartPar
Added support for plotting individual trends by group
\\
\sphinxhline
\sphinxAtStartPar
1.0.3
&
\sphinxAtStartPar
Bug correction
\\
\sphinxhline
\sphinxAtStartPar
1.0.2
&
\sphinxAtStartPar
Added support to parse full csv/excel automatically
\\
\sphinxhline
\sphinxAtStartPar
1.0.1
&
\sphinxAtStartPar
First implementation of algorithm
\\
\sphinxbottomrule
\end{tabulary}
\sphinxtableafterendhook\par
\sphinxattableend\end{savenotes}


\renewcommand{\indexname}{Python Module Index}
\begin{sphinxtheindex}
\let\bigletter\sphinxstyleindexlettergroup
\bigletter{c}
\item\relax\sphinxstyleindexentry{chowclassifier}\sphinxstyleindexpageref{source/chowclassifier:\detokenize{module-chowclassifier}}
\item\relax\sphinxstyleindexentry{chowclassifier.chow}\sphinxstyleindexpageref{source/chowclassifier:\detokenize{module-chowclassifier.chow}}
\end{sphinxtheindex}

\renewcommand{\indexname}{Index}
\printindex
\end{document}